\section{Three Gaps}
\label{sec:gaps}

Chapter~2 identified three specific things the published record leaves open. They are
stated here as gaps because each one is addressed by a design decision in Chapter~5.

\paragraph{Gap 1: no controlled architecture comparison.}
Studies comparing backbones on HAM10000 vary the split protocol, augmentation pipeline,
optimiser, batch size and epoch budget alongside the architecture. A reported difference
between two networks therefore confounds the network with the experimental setup. There
is no published result that isolates the architecture.

\paragraph{Gap 2: no factorial crossing architecture with imbalance strategy.}
Imbalance remedies are studied one at a time, usually against a single fixed backbone.
Whether the best remedy depends on which architecture it is applied to is an open
question, and one that matters: if the answer is yes, then any recommendation of a
remedy without reference to the backbone is incomplete.

\paragraph{Gap 3: aggregate metrics obscure the clinical picture.}
A model predicting nevus for every HAM10000 image achieves 66.9\% accuracy while
detecting no melanoma at all. Per-class F1 for the malignant categories is rarely the
headline anywhere in the literature, and averaged accuracy can improve while melanoma
detection degrades.

\section{Research Questions}
\label{sec:rqs}

Each question below has a pass condition or selection rule fixed before any experiment
was run. Fixing them in advance is what stops the analysis from drifting toward whatever
the data happens to support.

\subsection*{RQ1: the deployability gap}

\emph{Does the best MobileNetV2 configuration achieve a melanoma F1 within 0.05 of the
best ResNet-50 configuration?}

\begin{equation}
  \bigl| \mathrm{F1}_{\mathrm{MEL}}(\text{best MobileNetV2})
       - \mathrm{F1}_{\mathrm{MEL}}(\text{best ResNet-50}) \bigr| \le 0.05
  \label{eq:rq1}
\end{equation}

The threshold of 0.05 is a judgement, not a derivation, and is declared as such. It
represents the largest difference in melanoma F1 that would still leave the lightweight
model a defensible choice given roughly a tenfold reduction in size. A reviewer is
entitled to disagree with the number; what matters is that it was set before the results
existed rather than after.

``Best'' is defined by the RQ2 criterion below, applied within each architecture.
Section~\ref{sec:rq1-robustness} in Chapter~7 also reports the result under the
alternative reading in which ``best'' means highest macro F1, so that the conclusion can
be checked against both definitions.

EfficientNet-B3 does not compete for RQ1. It was added after the confirmatory design was
fixed, and allowing a later addition to displace one of the original arms would change
what the pre-registered question was asking.

\subsection*{RQ2: which imbalance strategy, for which architecture}

\emph{For each architecture, which of augmentation-only, weighted cross-entropy or
random oversampling maximises mean recall on the two clinically significant classes?}

\begin{equation}
  \text{best}(a) = \arg\max_{s \in S}
    \frac{\mathrm{recall}_{\mathrm{MEL}}(a, s) + \mathrm{recall}_{\mathrm{BCC}}(a, s)}{2}
  \label{eq:rq2}
\end{equation}

for architecture $a$ over strategies $S$.

Recall rather than F1 or accuracy, because the two error types are not equally costly
here. A melanoma classified as benign leaves untreated cancer. A benign lesion flagged
as suspicious costs a dermatologist's time and a patient's anxiety. In a screening
context that a specialist reviews downstream, the first error is much worse, and recall
is the metric that penalises it.

This choice has a consequence that Chapter~7 reports rather than hides: the criterion
can select a different configuration than macro F1 would. Both readings are given.

\subsection*{RQ3: what efficiency actually costs}

\emph{How do the architectures compare on trainable parameter count, serialised model
size, and CPU inference latency?}

No pass threshold: the question is descriptive. The measurement conditions are fixed
instead, and given in Chapter~6. Latency is measured on CPU because the deployment
scenario under discussion is an offline device without a GPU, and a GPU figure would not
speak to it.

\section{The EfficientNet-B3 Ablation}
\label{sec:ablation}

EfficientNet-B3 was added at the request of the thesis supervisor after the original six
runs were complete. It is handled as an ablation throughout, with two consequences.

First, it does not enter RQ1, as stated above. Second, its comparisons form a separate
statistical family with its own correction, described in Section~\ref{sec:stat-families}.
Pooling twenty-one comparisons into one family would tighten the correction applied to
the nine pre-registered ones, meaning a decision taken after seeing partial results
could change whether the original findings reach significance. Keeping the families
apart avoids that. It is a choice, and it is stated here rather than buried.

The ablation asks two things:

\begin{enumerate}
  \item At matched input resolution, does EfficientNet-B3 outperform either incumbent?
        Runs E7 to E9 use 224 pixels, identical to the other six, so any difference is
        attributable to the architecture.
  \item Does B3's native 300-pixel input account for its advantage, where one exists?
        Runs E10 to E12 repeat E7 to E9 at 300 pixels with everything else unchanged,
        making each pair a controlled test of resolution alone.
\end{enumerate}

The second question is the one the literature cannot currently answer, for the reason
given in Section~\ref{sec:lit-arch}: published comparisons give each architecture its
own native resolution and therefore vary two factors simultaneously.

\section{Scope Exclusions}
\label{sec:scope}

The following are excluded deliberately. Each is a real limitation and is revisited in
Chapter~9.

\paragraph{Focal loss.} Excluded to keep the factorial at twelve runs rather than
sixteen. Cited in Chapter~2 as the method used by the closest published
comparator~\cite{le2020classweightedfocal}.

\paragraph{Additional architectures.} DenseNet, ConvNeXt and Vision Transformers appear
in the surveyed corpus~\cite{wei2023densenetconvnext} and are not included. The question
being answered is what a deployable network costs against a standard baseline, and each
extra backbone dilutes that without sharpening it.

\paragraph{Ensembles and fusion.} Several surveyed studies report gains from
ensembling~\cite{chiu2025ensemble, alruwaili2025fusion}. An ensemble of $n$ models
occupies $n$ times the space, which is the opposite of the deployment constraint
motivating this work.

\paragraph{External validation.} All runs use HAM10000 only. Cross-dataset
generalisation to ISIC 2019 or 2020 would require a second data pipeline and would
reintroduce exactly the protocol variation that Gap~1 is about.

\paragraph{Explainability.} Grad-CAM~\cite{selvaraju2017gradcam} and related methods are
not applied. This is a limitation rather than a considered exclusion: attention maps
would help confirm the models attend to lesions rather than to acquisition artefacts.

\paragraph{Deployment itself.} The thesis quantifies the trade-off. It does not build a
mobile application, and reports no on-device measurements.

\section{Objectives}

\begin{enumerate}
  \item Implement a single training pipeline in which architecture, imbalance strategy
        and input resolution are the only variables, with everything else held fixed and
        verifiable.
  \item Execute the $2\times3$ confirmatory factorial and the six-run EfficientNet-B3
        ablation on identical hardware from one shared data split.
  \item Report per-class precision, recall and F1 for all seven classes in every
        configuration.
  \item Answer RQ1 and RQ2 against the criteria fixed in Equations~\ref{eq:rq1}
        and~\ref{eq:rq2}.
  \item Measure RQ3 quantities directly rather than quoting them from the architecture
        papers.
  \item Test every declared comparison with McNemar's test under Holm-Bonferroni
        correction, in the two families described above.
  \item Persist per-run configuration, training history, test metrics and per-image
        predictions so that every number in this document is regenerable from stored
        artefacts.
\end{enumerate}
